\documentclass[../../../include/open-logic-section]{subfiles}

\begin{document}

\olfileid{sth}{replacement}{intro}
\olsection{Pambuka}

Panggantian iku skema aksioma kang mbedakake
$\ZF$ lan~$\Z$. Kita wis nggunakake skema iki
kanthi bebas ing saindenging
\crefrange{sth:ordinals::chap}{sth:spine::chap}.
Ing bab iki, pungkasane kita bakal ngrembug
pitakonan: apa Panggantian bisa diwenehi dhasar?

Supaya pitakonan iki luwih cetha, elinga yen
Panggantian pancen cukup \emph{kuwat}. Ing bab
iki (lan maneh ing
\olref[sth][card-arithmetic][fix]{sec}),
kita bakal weruh sepira kuwate. Nanging iki
uga nuduhake yen Panggantian pancen perlu
diwenehi dhasar.

Kita bakal ngrembug rong jinis pambeneran. Kanthi
garis gedhe, pambeneran \emph{ekstrinsik} nyoba
mbenerake aksioma saka asil kang diasilake;
pambeneran \emph{intrinsik} nyoba mbenerake
aksioma kanthi nuduhake yen aksioma iku didhukung
dening konsep-konsep matematika kang dirembug.
Ing bab iki, kita bakal luwih ngerti tegesé
prabedan iki, nanging sing dirembug mung
saperangan cilik saka perkara kang luwih amba.
Kanggo katrangan luwih jangkep, delengen utamane
\citeauthor{Maddy1988a} (\citeyear{Maddy1988a}
lan \citeyear{Maddy1988b}).

\end{document}
